\documentclass[10pt,a4paper]{amsart}
\usepackage[margin=19mm]{geometry}
\usepackage{amsmath,amssymb,mathtools}
\usepackage[scr=rsfs,cal=euler]{mathalfa}
\usepackage{xfrac}
\usepackage[unicode,hidelinks,bookmarks=false]{hyperref}
\newcommand{\ie}{i.e.\ }
\newcommand{\cf}{cf.\ }
\newcommand{\resp}{respectively\ }
\newcommand{\Subsection}{\subsection}
\providecommand{\nobreakdash}{\nobreak\hbox{-}}
\setlength{\emergencystretch}{2em}
\multlinegap0pt
\usepackage{csquotes}
\usepackage{mathtools}
\usepackage{amssymb}
\usepackage{xcolor}
\newtheorem{theorem}{Theorem}
\def\gcd{{\rm gcd}}
\def\ord{{\rm ord}}
\newcommand{\res}{\operatorname{Res}}
\title{Resultants of dynatomic polynomials of $x^d$}
\author{Chih-Chiang Kao}
\date{Source: arXiv:2608.27016v1 (CC BY 4.0)}
\begin{document}
\maketitle

\noindent\emph{Source and licence.} Resultants of dynatomic polynomials of $x^d$. Version arXiv:2608.27016v1 (CC BY 4.0), distributed by arXiv under the Creative Commons Attribution 4.0 International licence, http://creativecommons.org/licenses/by/4.0/, as recorded in the arXiv licence metadata for this exact version and shown on the abstract page https://arxiv.org/abs/2608.27016v1. Full licence notice: \texttt{licenses/arxiv-2608.27016-CC-BY-4.0.txt}.\par\smallskip

\noindent\textit{Editorial scope.} Counted unit: the article's theorem theorem (source0.tex lines 542--549). The article's complete original proof of it is delivered with it (source0.tex lines 551--584, one \texttt{proof} environment of 34 lines). No mathematics is written, repaired or re-derived: the statement and the proof are the article's own lines, in the article's own notation, and no retained line is substituted or re-flowed.  No further source-local context is needed: every cross-reference of every retained line resolves inside the two delivered spans.  Nothing was omitted from any retained span, so the package carries no omission record.  No retained line contains a citation, so the wrapper carries no bibliography and no bibliography entry.  No retained line contains a figure environment, an image-inclusion command or drawing code, so no image is delivered and the package declares no asset dependency.  Editorial formatting only, in addition: an AMS \texttt{amsart} preamble that restates the article's own declarations and macros used by the retained lines, namely the environments \texttt{theorem} on separate counters, together with the standard packages those lines need.  The retained environments print in this document as Theorem 1 in order of appearance, whereas the article prints the same items with the numbering of its own full document.  The complete native source of the article as distributed in the arXiv e-print for this exact version is bundled unchanged as \texttt{sources/208677/source0.tex}.
\begin{theorem}
    Let $1<n<m$. The following statements hold. 
    \begin{enumerate}
        \item $\res(\Phi_{d,n},\Phi_{d,m}) = 1$ if and only if $n\nmid m$.
        \item Let $p$ be a prime divisor of $\Phi_{m}(d)$. If $n\mid m$, then
        $$v_p(\res(\Phi_{d,n},\Phi_{d,m}))\geq v_p(\Phi_m(d))\left(\sum_{k\mid n}\mu\left(\frac{n}{k}\right)d^k\right).$$
    \end{enumerate}
\end{theorem}

\begin{proof}
    \begin{enumerate}
        \item First, we suppose $n\nmid m$. Since $n\nmid m$ and $n<m$, it follows that $a\nmid b$ and $b\nmid a$ by Lemma 4.1 and Corollary 4.2. Therefore, $\res(\Phi_a,\Phi_b) = 1$ for every pair $(a,b)$. Hence, $\res(\Phi_{d,n},\Phi_{d,m}) = 1$.

        The converse direction follows immediately from Theorem 4.4(2). For any prime divisor $p$ of $\Phi_m(d)$, since $\sum_{k\mid n}\mu(n/k)d^k = \deg\Phi_{d,n}$, it follows that
        $$v_p(\res(\Phi_{d,n},\Phi_{d,m})) \geq v_p(\Phi_m(d))\cdot\sum_{k\mid n}\mu(n/k)d^k\geq 1.$$
        Hence, if $n\mid m$, then $\res(\Phi_{d,n},\Phi_{d,m})>1$.
        \item Let $p$ be a prime divisor of $\Phi_m(d)$. Note that
        $$\res(\Phi_{d,n},\Phi_{d,m}) = \prod_{\ord_a(d) = n}\prod_{\ord_b(d) = m}\res(\Phi_a,\Phi_b).$$
        By Theorem 2.2, $\res(\Phi_a,\Phi_b) = p^{\varphi(a)}$ if and only if $b/a$ is a power of $p$. Hence,
        \begin{align}
            v_p(\res(\Phi_{d,n},\Phi_{d,m})) &= \sum_{\ord_{a}(d) = n}\sum_{\ord_b(d) = m}v_p(\res(\Phi_a,\Phi_b))\notag\\
            &= \sum_{\ord_a(d) = n}\varphi(a)\cdot |\{t\geq 1: \ord_{ap^t}(d) = m\}|.
        \end{align}
        Define $T(a,p,d,m) = \{t\geq 1:\ord_{a\cdot p^t}(d) = m\}$. Recall that $S(p,d,m) = \{t\geq 1:\ord_{p^t}(d) = m\}$. We claim that
        $$|T(a,p,d,m)|\geq |S(p,d,m)|$$
        whenever $n\mid m$ and $\ord_a(d) = n$. Since $p\mid \Phi_m(d)$, it follows that $p\mid d^m-1$. We denote $v_p(d^m-1)$ by $r$. For each prime divisor $q$ of $m$, we denote $v_p(d^{m/q}-1)$ by $r_q$. Then $t\in S(p,d,m)$ if and only if $p^t\mid d^m-1$ and $p^t\nmid d^{m/q}-1$ for all prime divisors $q$ of $m$. Hence, $\max_{q\mid m}r_q<t\leq r$. Thus, $S(p,d,m) = \{t\geq 1:\max_{q\mid m}r_q<t\leq r\}$ and so $|S(p,d,m)| = r-\max_{q\mid m}r_q$.

        Write $a = a'p^k$, where $k = v_p(a)$ and $\gcd(a',p) = 1$. Then $a\cdot p^t = a'\cdot p^{t+k}$ and 
        $t\in T(a,p,d,m)$ if and only if $a'\cdot p^{t+k}\mid d^m-1$ and $a'\cdot p^{t+k}\nmid d^{m/q}-1$ for all prime divisors $q$ of $m$. To simplify the notation, define
        $$Q(n,m) = \{\text{prime divisors }q\text{ of }m:n\text{ divides }m/q\}.$$
        Suppose $q\in Q(n,m)$. Since $a'\mid a$, it follows that $a'\mid d^n-1$ and so $a'\mid d^{m/q}-1$. In this case, $t+k\leq r$ and $r_q<t+k$. Hence, $T(a,p,d,m)\supseteq \{t\geq 1:\max_{q\in Q(n,m)}(r_q-k)<t\leq r-k\}$, and
        $$|T(a,p,d,m)|\geq r-k-\max_{q\in Q(n,m)}(r_q-k) = r-\max_{q\in Q(n,m)}r_q.$$
        Since $Q(n,m)$ is a subset of the set of prime divisors $q$ of $m$, one has $\max_{q\in Q(n,m)}r_q\leq \max_{q\mid m}r_q$. Hence, 
        $$|T(a,p,d,m)| \geq r-\max_{q\in Q(n,m)}r_q\geq r-\max_{q\mid m}r_q = |S(p,d,m)|.$$
        Substituting this estimate into the equation (4.1) yields
        \begin{align*}
            v_p(\res(\Phi_{d,n},\Phi_{d,m})) &= \sum_{\ord_a(d) = n}\varphi(a)\cdot|T(a,p,d,m)|\\
            &\geq \sum_{\ord_a(d) = n}\varphi(a)\cdot|S(p,d,m)|\\
            &= v_p(\Phi_{m}(d))\cdot\sum_{\ord_a(d) = n}\varphi(a)&&\text{by Lemma 3.9}\\
            &= v_p(\Phi_m(d))\sum_{k\mid n}\mu(n/k)d^k&&\text{by Proposition 3.3(1)}.
        \end{align*}
    \end{enumerate}
\end{proof}

\end{document}
